% SEGMENT OLP-0616-S01
% Part: methods
% Chapter: induction
% Section: strong-induction

\documentclass[../../../include/open-logic-section]{subfiles}

\begin{document}

\olfileid{mth}{ind}{str}

\olsection{வலுவான தொகுத்தறிதல்}

மேலே விவாதித்த தொகுத்தறிதல் கொள்கையில் $P(0)$
என்பதையும், $P(k)$ எனில் $P(k+1)$ என்பதையும்
நிறுவுகிறோம். இரண்டாம் பகுதியில் $P(k)$ மெய்
என்று எடுத்துக்கொண்டு, அந்த எடுகோளைப் பயன்படுத்தி
$P(k+1)$ ஐ நிறுவுகிறோம். இதற்குச் சமமாக,
$P(k-1)$ ஐ எடுத்துக்கொண்டு $P(k)$ ஐ
நிறுவலாம். எந்த எண்ணிலிருந்தும் அதற்கு அடுத்த
எண்ணுக்கு உய்த்தறிய முடிவதே முக்கியம்: முன்னுள்ள
எண்ணுக்கு ஒரு கூற்று மெய் என்ற எடுகோளின்கீழ்,
அந்தக் கூற்றை எந்த எண்ணுக்கும் நிறுவ முடிவது.

% SEGMENT OLP-0616-S02
தொகுத்தறிதல் கொள்கையின் வேறொரு வடிவத்தில்,
$k$ க்கு முந்திய $k-1$ என்ற எண்ணுக்கு மட்டுமே
கூற்று மெய் என்று எடுத்துக்கொள்ளாமல்,
~$k$ ஐவிடச் சிறிய எல்லா எண்களுக்கும் மெய்
என்று எடுத்துக்கொண்டு, அதன் மூலம்~$k$ க்கான
கூற்றை நிறுவுகிறோம். இதுவும் எல்லா~$n \in \Nat$
க்கும் $P(n)$ ஐத் தருகிறது. ஏனெனில் $P(0)$
ஐ நிறுவியதும்,~$1$ ஐவிடச் சிறிய எல்லா
எண்களுக்கும் $P$ மெய் என்று நிறுவியிருக்கிறோம்.
எல்லா $l<k$ க்கும் $P(l)$ எனில் $P(k)$
என்று தெரிந்தால், குறிப்பாக $k=1$ க்கும்
இது பொருந்தும். ஆகவே $P(1)$ என்று முடிவு
செய்யலாம். இப்போது $P(0)$, $P(1)$ இரண்டும்,
அதாவது எல்லா $l<2$ க்கும் $P(l)$,
நிறுவப்பட்டுள்ளன. மேலும் எல்லா $l<2$ க்கும்
$P(l)$ எனில் $P(2)$ என்ற நிபந்தனையும்
நமக்குள்ளது; ஆகவே~$P(2)$ என்று முடிவு
செய்யலாம். இப்படியே தொடரலாம்.

% SEGMENT OLP-0616-S03
உண்மையில் ``எல்லா $k$ க்கும், எல்லா $l<k$
க்கும் $P(l)$ எனில் $P(k)$'' என்ற பொதுவான
நிபந்தனைக் கூற்றை நிறுவ முடிந்தால், $P(0)$
ஐத் தனியாக நிறுவத் தேவையில்லை; அதிலிருந்தே
அது கிடைக்கிறது. ஏனெனில் ``எல்லா $l<k$
க்கும் $P(l)$'' போன்ற பொதுக் கூற்று,
$l<k$ ஆகும் எண் எதுவும் இல்லாவிட்டாலும்
மெய். இது வெற்றிடக் கணக்கீட்டின் ஒரு நிகழ்வு:
``எல்லா $A$ களும் $B$ களே'' என்பது
$A$ எதுவும் இல்லாதபோதும் மெய்;
$\lforall[x][(!A(x) \lif !B(x))]$ என்பது
~$!A(x)$ ஐ நிறைவேற்றும் $x$ எதுவும்
இல்லாவிட்டால் மெய். இங்கே அதன் முறைப்படியான
வடிவம் ``$\lforall[l][(l < k \lif P(l))]$'';
$l < k$ ஆகும் எண் எதுவும் இல்லாவிட்டாலும்
இது மெய். $k=0$ எனில் அப்படித்தான்:
$l<0$ ஆகும் எண் இல்லை. ஆகவே $P$ எதுவாக
இருந்தாலும் ``எல்லா $l<0$ க்கும் $P(0)$''
மெய். இதனால் ``எல்லா $l<k$ க்கும்
$P(l)$ எனில் $P(k)$'' என்பதற்கான நிறுவலே
தானாக~$P(0)$ ஐயும் நிறுவுகிறது.

% SEGMENT OLP-0616-S04
~$k$ க்கான கூற்றை நிறுவ $k-1$ க்கான
கூற்றை மட்டும் நம்ப முடியாமல், ஒன்று அல்லது
பல $l<k$ களுக்கான கூற்று மெய் என்ற
எடுகோள் தேவைப்படும் போது இந்த வடிவம்
பயனுள்ளதாகும்.

\end{document}
